% why socially intelligent agents are important
Recent advances in human-esque collaborative Large Language Model (LLM)-based agents~\citep{zhang2023exploring} have raised new questions about the feasibility of modeling human social learning with socially intelligent agent~\citep{socially-intelligent-machines}. 
% why directly using LLMs for social tasks is not desired
Existing work has explored LLMs' social capabilities in one-shot instruction following~\citep{sap2023neural, gandhi2023understanding, shapira2023clever} settings. While showing impressive social skills, standard LLMs struggle with navigating social situations in goal-oriented social interactions~\citep{sotopia}. 
% what are existing methods of training/enabling better social agents
Current attempts to construct socially intelligent agents have focused on engineering prompts to influence LLMs' actions~\citep{park2022socialsimulacra, fischer2023reflective} and on fine-tuning LLM agents with LLM-generated datasets~\citep{liu2023training}. 
% what are the limitations of existing methods
However, prompt engineering-based methods are difficult to generalize to diverse social scenarios and existing fine-tuning methods improve agents' social alignment only under single-turn instruction-response settings. 
There is a lack of research in training LLM agents to acquire social intelligence which applies to multi-turn goal-oriented social interactions. 


% What do we propose
In this work, we propose the first social interaction learning framework for training socially intelligent LLM agents that can navigate multi-turn goal-oriented social interactions. 
% From where we draw inspirations
This framework draws inspirations from inferential social learning in cognitive science, where humans learn not by simply imitating, but by deciding when to imitate depending on their sophisticated inferences of who is doing what and why~\citep{inferential-social-learning}, which is \textit{learning from social interaction}~\citep{premack_woodruff_1978, Tomasello2021becoming}.
The two stages of the framework: supervised fine-tuning and offline reinforced self-training, correspond to humans' social learning through social imitation and interaction. 
% why our approach solves the problem
The first stage warms up LLM agents by fine-tuning on LLM-generated social interaction data. 
The second stage enables LLM agents to autonomously expand social scenarios and self-generate social interaction data for future steps of fine-tuning, with reward signals purely from LLMs. 


% What are our results
We train and evaluate LLM agents in the SOTOPIA interactive simulated social environment~\citep{sotopia}. 
The results show that our framework improves LLM agents' social performance in multi-turn goal-oriented social interactions. 
And the social performance of the trained LLMs surpasses their instruction fine-tuned counterparts, indicating the potential of our training framework to become an alternative to instruction fine-tuning in improving social intelligence. 
In addition, our approach enables LLM agents to align with human values and social norms by simulating morally salient social scenarios~\citep{hendrycks2021what}. We evaluate the agents on metrics of social norms (e.g., safety, toxicity, the ability to keep secrets) as the method for human values verification~\cite{ji2024ai}. In a broader context, our paper aims to make progress toward an open question in the LLM community: can LLM agents with higher social intelligence achieve better human values alignment?


% Our contributions are as follows: 1) We propose the first framework for training socially intelligent LLM agents that can navigate multi-turn goal-oriented social interactions. Our experiments show an increase in LLM agents' social intelligence and their alignment with human values after training with our framework. 2) We create and publish the trained LLMs with higher social intelligence and better human values alignment. We also provide a comprehensive empirical analysis of our proposed framework and the trained LLMs.